The Reflexion loop is the outer self-critique layer of LRAM. It treats every rejected trajectory as a learning signal and writes a natural-language reflection to episodic memory. Subsequent inference and training steps read from this memory to bias the policy away from known failure modes.

\subsection{Episodic Memory}
The episodic memory is a bounded buffer of textual reflections, each indexed by the failure mode it addresses. Memory entries have decay weights that decrease with age, preventing the policy from overfitting to outdated critiques.

\subsection{Reflection Policy}
A reflection is generated by a dedicated teacher call triggered whenever the tier-3 validator rejects a candidate. The teacher is asked to produce a concise natural-language explanation of the rejection reason and a suggested prior to bias subsequent attempts. The explanation is stored verbatim and the suggested prior is embedded numerically for use as a soft constraint in the next Prajna compression.

\subsection{Convergence Condition}
A Reflexion loop converges when the rejection rate drops below a target threshold for a sustained number of training steps. The formal convergence condition is stated as an equation in Section 5 and Appendix equation set. In practice we monitor three diagnostic time series: the rejection rate, the Reflexion buffer entropy, and the circuit-attribution stability.

\subsection{Interaction with Teacher-Student Distillation}
The Reflexion loop and the teacher-student distillation loop operate at different timescales. The distillation loop is gradient-based and slow; the Reflexion loop is text-based and fast. The two are composed by having Reflexion entries contribute a soft penalty term to the distillation loss, so that textual critiques influence gradient updates without the engineering burden of reward modelling.

\subsection{Diagnostic Metrics}
For each Reflexion epoch we log the number of new entries, the number of retired entries, the mean decay weight, and the correlation between Reflexion entries and subsequent tier-3 rejection reductions. These metrics are reported in Section 12.
